\exercisesection{4}

\begin{enumerate}
\def\labelenumi{\arabic{enumi}.}
\item
  设 \((\mathbf{X}_{\alpha}, f_{\alpha\beta})\) 是拓扑空间的一个逆向系，其指标集是有向的，\(X = \varprojlim X_{\alpha}\) ，而 \(f_{\alpha}\) 是 X 到 \(X_{\alpha}\) 的典范映射。假设对所有 \(\alpha, f_{\alpha}\) 都是满射；证明：若诸 \(X_{\alpha}\) 都不可约，则 X 不可约（参看《一般拓扑学》第 I 章 §4 第 4 节命题 9 的推论）。特别地，不可约空间的任意乘积都不可约。
\item
  不可约空间中的一点 \(x\) 称为该空间 \(\mathbf{X}\) 的泛点，如果 \(\{x\}\) 在 \(\mathbf{X}\) 中处处稠密。
\end{enumerate}

\begin{enumerate}
\def\labelenumi{(\alph{enumi})}
\tightlist
\item
  若 \(\mathbf{X}\) 是一个 Kolmogoroff 空间（《一般拓扑学》第 I 章 § 1 习题 2），它至少有一个泛点；若 \(\mathbf{X}\) 是一个可达空间（《一般拓扑学
\end{enumerate}

》第 I 章 § 8 习题 1），则除非它仅由一点组成，它不可能有泛点。

\begin{enumerate}
\def\labelenumi{(\alph{enumi})}
\setcounter{enumi}{1}
\item
  给出一个无限的不可约可达空间的例子（参看《一般拓扑学》第 I 章 § 8 习题 5）。
\item
  设 X、Y 是两个都有泛点的不可约空间；再假设 Y 只有唯一的泛点 y。设 \(f: X \to Y\) 是一个连续映射；为使 \(f(X)\) 在 Y 中稠密，必须且只需对 X 的每个泛点 x 都有 \(f(x) = y\) 。
\item
  设 \((\mathbf{X}_{\alpha}, f_{\alpha \beta})\) 是不可约空间的一个逆向系，其指标集是有向的；假设每个 \(\mathbf{X}_{\alpha}\) 都有唯一的泛点 \(x_{\alpha}\) 。证明：若对 \(\alpha \leqslant \beta, f_{\alpha \beta}(\mathbf{X}_{\beta})\) 在 \(\mathbf{X}_{\alpha}\) 中稠密，则 \(\mathbf{X} = \varprojlim \mathbf{X}_{\alpha}\) 不可约且只有唯一的泛点（方法同习题 1）。
\end{enumerate}

\begin{enumerate}
\def\labelenumi{\arabic{enumi}.}
\setcounter{enumi}{2}
\item
  设 X 是一个拓扑空间，\((Y_{\alpha})\) 是 X 的子空间的一个右有向族；证明：若每个子空间 \(Y_{\alpha}\) 都不可约，且 X 是族 \((Y_{\alpha})\) 之并，则 X 不可约。由此给出一个例子：X 是 Kolmogoroff 空间，每个 \(Y_{\alpha}\) 都有泛点，但 X 没有泛点。
\item
  设 Y 是只有唯一泛点 y 的不可约空间，X 是一个拓扑空间，\(f: X \rightarrow Y\) 是一个连续映射。
\end{enumerate}

\begin{enumerate}
\def\labelenumi{(\alph{enumi})}
\item
  设 \(\mathbf{Z}\) 是 \(\mathbf{X}\) 的一个不可约分支，它与 \(f^{-1}(y), f(\mathbf{Z})\) 中的前者相交；则其像在 \(\mathbf{Y}\) 中稠密。
\item
  给出一个例子，其中 \(f(\mathbf{X})\) 在 \(\mathbf{Y}\) 中稠密，但 \(f^{-1}(y)\) 为空（取 \(\mathbf{X}\) 为 \(\mathbf{Y}\) 的一个子空间）。
\item
  证明：若 \(\mathbf{Z}\) 有一个泛点 \(z\) 且 \(f(\mathbf{Z})\) 在 \(\mathbf{Y}\) 中稠密，则 \(f(z) = y\) ，且 \(z\) 是 \(\mathbf{Z} \cap f^{-1}(y)\) 的一个泛点。
\end{enumerate}

\begin{enumerate}
\def\labelenumi{\arabic{enumi}.}
\setcounter{enumi}{4}
\item
  设 G 是一个连通的半拓扑群（《一般拓扑学》第 III 章 § 1 习题 2）。证明：若 G 只有有限多个不可约分支，则它必然不可约（注意诸不可约分支都可由包含单位元 e 的那个不可约分支经左平移或右平移得到）。
\item
  设 \(\mathbf{X}\) 是一个拓扑空间。
\end{enumerate}

\begin{enumerate}
\def\labelenumi{(\alph{enumi})}
\item
  设 x 是 X 的一点，U 是 x 的一个只有有限多个不可约分支的开邻域。证明存在 x 的一个邻域 V，使得包含于 V 的每个邻域都是连通的。
\item
  假设 X 的每一点都有一个只含有限多个不可约分支的开邻域。证明下列性质等价：
\end{enumerate}

(α) X 的诸不可约分支都是开的。

(β) X 的诸不可约分支与它的诸连通分支相同。

(γ) X 的诸连通分支都不可约。

(δ) X 的两个互不相同的不可约分支不相交。

\begin{enumerate}
\def\labelenumi{\arabic{enumi}.}
\setcounter{enumi}{6}
\item
  设 X 是一个紧可度量化空间，R 是 X 上的一个开等价关系，使得商空间 X/R 不可约。证明存在一点 \(x \in X\) ，它模 R 的类处处稠密。（首先注意 X 的每个关于 R 饱和的开子集都处处稠密。然后利用 Baire 定理。）
\item
  证明两个 Noether 空间的乘积是 Noether 空间。（若 X 与 Y 是 Noether 空间，A 是 \(X \times Y\) 的一个开子集，证明：对所有 \(x \in pr_{1}A\) ，存在 x 的一个开邻域 V，使得 \(V \times A(x) \subset A\) ，并由此推出 A 是拟紧的。）
\item
  \begin{enumerate}
  \def\labelenumii{(\alph{enumii})}
  \tightlist
  \item
    证明交换环的素谱是一个 Kolmogoroff 空间（《一般拓扑学》第 I 章 § 1 习题 2），其中每个不可约分支都有（唯一的）泛点。
  \end{enumerate}
\end{enumerate}

\begin{enumerate}
\def\labelenumi{(\alph{enumi})}
\setcounter{enumi}{1}
\tightlist
\item
  设 A 是一个整环，\(\xi = (0)\) 是 \(\mathbf{X} = \operatorname{Spec}(\mathbf{A})\) 的泛点。为使点 \(\xi\) 在 X 中孤立，必须且只需 A 的诸 \(\neq 0\) 素理想的交是一个 \(\neq 0\) 的理想。给出一个具有此性质的局部环的例子。
\end{enumerate}

\begin{enumerate}
\def\labelenumi{\arabic{enumi}.}
\setcounter{enumi}{9}
\item
  设 A 是一个环，N 是它的幂零根。为使 A 的谱是离散的，必须且只需 A/N 是有限多个域的直复合。
\item
  设 K 是一个域，A 是多项式环 K{[}X, Y{]}，B 是商环 A/(XY)；证明 Spec(B) 是连通的，但有两个互不相同的不可约分支。
\item
  设 Y 是一个完全正则空间，\(\mathbf{A} = \mathcal{C}^{\infty}(\mathbf{Y}; \mathbf{R})\) ，\(\mathbf{B} = \mathcal{C}(\mathbf{Y}; \mathbf{R})\) 。证明 Y 的 Stone-Čech 紧化 \(\beta Y\) 典范地等同于 \(\operatorname{Spec}(A)\) 的一个处处稠密子空间，并且典范地等同于 \(\operatorname{Spec}(B)\) 的一个处处稠密子空间。
\item
  设 \(\mathbf{A}\) 是一个环，\(\mathbf{X} = \operatorname{Spec}(\mathbf{A})\) 是它的谱。
\end{enumerate}

\begin{enumerate}
\def\labelenumi{(\alph{enumi})}
\item
  若 X 的一个子集 F 是闭的，它就具有下列两个性质：(α) 对所有 p ∈ F，有 V(p) ⊂ F；(β) 对所有 p ∉ F，存在 V(p) 的一个闭子集，它包含 F ∩ V(p) 但不包含 p。
\item
  假设 X 是 Noether 的。证明 X 的每个具有 (a) 中性质 \((\alpha)\) 与 \((\beta)\) 的子集 F 在 X 中是闭的（考虑 \(\overline{F}\) 的不可约分支，并利用习题 9 (a)）。
\item
  沿用习题 12 的记号，取 Y 为 R 的区间 {[}0, 1{]}。证明 Y 在 X = Spec(A) 中不是闭的（参看 § 2 习题 15），但满足 (a) 的条件 ( \(\alpha\) ) 与 ( \(\beta\) )。
\end{enumerate}

\begin{enumerate}
\def\labelenumi{\arabic{enumi}.}
\setcounter{enumi}{13}
\tightlist
\item
  设 A 是一个环，\(\mathbf{X} = \operatorname{Spec}(\mathbf{A})\) 是它的谱，P 是 A 的极小素理想组成的集合。
\end{enumerate}

\begin{enumerate}
\def\labelenumi{(\alph{enumi})}
\item
  设 \(R_{\mathfrak{p},\mathfrak{p}^{\prime}}\) 是对称且自反的关系：“存在 A 的一个素理想 \(\mathfrak{p}''\) ，它包含 \(\mathfrak{p} + \mathfrak{p}'''\) ，并且位于 P 的元素 \(\mathfrak{p},\mathfrak{p}'\) 之间”；设 S 是这样的等价关系：它的图是包含 R 的图的诸等价关系的图中最小的一个（《集合论》第 II 章 § 6 习题 10）。证明：若 I 是关于 S 的一个等价类，则集合 \(V_{I} = \bigcup_{\mathfrak{p}\in I}V(\mathfrak{p})\) 是连通的。
\item
  证明：若 P 是有限的（特别地，若 A 是 Noether 的），则诸 \(V_{I}\) 是 X 的连通分支。这一结果能否推广到 P 无限的情形（参看习题 12）？
\end{enumerate}

¶ 15. 设 A 是一个环，a 是 A 的一个有限生成理想。证明下列性质等价：(α) \(a^{2} = a\) ；(β) A 由一个幂等元生成；(γ) V(a) 在 X = Spec(A) 中既开又闭，且 a 在满足 r(b) = r(a) 的诸理想 b 中是最小的（利用 § 2 命题 4 的推论 3）。给出一个例子，其中 \(a^{2} = a\) ，a 不是有限生成的，且不含幂等元（参看 § 2 习题 15）。

¶ 16. (a) 设 A 是一个绝对平坦的（交换）环（第 I 章 §2 习题 17）。证明 X = Spec(A) 是一个全不连通的紧空间，A 的每个素理想 p 都是极大的，并且 \(A_{p}\) 典范地同构于域 A/p（首先证明：对所有 f ∈ A，V(f) 在 X 中既开又闭；另一方面利用 §3 习题 9）。

\begin{enumerate}
\def\labelenumi{(\alph{enumi})}
\setcounter{enumi}{1}
\tightlist
\item
  证明映射 \(\mathfrak{a}\mapsto\mathrm{V}(\mathfrak{a})\) 是 A 的理想组成的集合到 X 的闭子集组成的集合上的一个双射。此外，下列条件等价：( \(\alpha\) ) V(a) 是开的；( \(\beta\) ) a 是有限生成的；( \(\gamma\) ) a 由一个幂等元生成；( \(\delta\) ) A/a 是一个投射 A-模。
\end{enumerate}

* (c) 设 \(\tilde{A}\) 是 \(X\) 上的结构环层。证明每个 \(\tilde{A}\) -模 \(\mathcal{F}\) 都形如 \(\tilde{E}\) ，其中 \(E\) 是一个在相差一个唯一同构的意义下确定的 A-模（取 \(E = \Gamma(X, \mathcal{F})\) ，并利用 \(X\) 的每一点都有一个由既开又闭的邻域组成的基本系这一事实）。设 \(a\) 是 \(A\) 的任意一个理想，证明 \(\operatorname{Hom}_A(a, E)\) 等同于 \(\Gamma(X - V(a), \tilde{E})\) （注意 \(X - V(a)\) 是诸 \(X_e\) 之并，其中 \(e\) 取遍幂等元 \(e \in a\) 组成的集合）。由此推出：为使 \(E\) 是一个内射 A-模，必须且只需 \(\tilde{E}\) 是松弛的。*

\begin{enumerate}
\def\labelenumi{(\alph{enumi})}
\setcounter{enumi}{3}
\tightlist
\item
  设 A 是一个环，\(\mathfrak{N}\) 是它的幂零根。下列性质等价：
\end{enumerate}

(α) \(A/\mathfrak{N}\) 是绝对平坦的。

(β) \(\mathbf{X} = \operatorname{Spec}(\mathbf{A})\) 是 Hausdorff 的。

(γ) X 的每一点都是闭的（换言之，A 的每个素理想都是极大的）。

（利用 § 3 的习题 9。）

¶ 17. * (a) 设 X 是一个全不连通的紧空间，O 是 X 上的一个环层，使得对所有 \(x \in X\) ，\(\mathcal{O}_x\) 都是一个交换域。证明：若 A = \(\Gamma(X, \mathcal{O})\) ，则环 A 是绝对平坦的，且它的谱（连带环层 \(\tilde{A}\) ）典范地等同于 X（连带环层 O）。（注意：对所有 \(f \in A\) ，诸 \(x \in X\) （满足 \(f_x = 0\) 者）在 X 中组成一个既开又闭的集合，并由此推出存在 \(g \in A\) 使 \(f = gf^2\) 。）*

\begin{enumerate}
\def\labelenumi{(\alph{enumi})}
\setcounter{enumi}{1}
\item
  设 X 是一个全不连通的紧空间，K 是一个（交换）域，A 是 X 到 K 的局部常值映射组成的环。证明 A 是绝对平坦的（直接论证或利用 (a)）。
\item
  设 A 是定义在 I = {[}0, 1{[} 上的实值阶梯函数组成的环，这些函数的不连续点都形如 \(k/2^n\) （ \(n \in \mathbf{N}, 0 \leqslant k < 2^n\) ），并且都是右连续的。设 \(\mathfrak{S}\) 是诸半开区间 \([x, y[ \subset I\) （其端点形如 \(k/2^n\) ）的有限并组成的集合。对每个 \(f \in A\) ，
\end{enumerate}

集合 \(Z(f) = f^{-1}(0)\) 属于 \(\mathfrak{S}\) 。若 a 是 A 的一个理想，则集合 \(\mathfrak{B}\) 由诸 \(Z(f)\) （其中 f 取遍 a）组成，且满足：(1) \(\mathfrak{B}\) 中两个集合的交属于 \(\mathfrak{B}\) ；(2) \(\mathfrak{S}\) 中每个包含 \(\mathfrak{B}\) 中一个集合的集合都属于 \(\mathfrak{B}\) 。证明其逆。为使 a 是素的，必须且只需 \(\mathfrak{B}\) 是收敛到 {[}0, 1{]} 中一点的一个滤子的基。证明 A 是一个绝对平坦的约化环，并且对 A 的每个极大理想 m，\(A_{m}\) 都同构于域 R。描述空间 Spec(A)。

¶ 18. (a) 设 Y 是一个完全正则空间。证明下列条件等价：

\((\alpha)\) \(\mathcal{C}(\mathbf{Y};\mathbf{R})\) 的每个素理想都是极大的。

(β) Y 的开子集的每个可数交都是开的。

(γ) Y 上的每个实值连续函数都是局部常值的。

(δ) 环 \(\mathcal{C}(\mathbf{Y};\mathbf{R})\) 是绝对平坦的。

（利用 § 2 习题 15 (f)。）

若这些条件成立，则空间 \(\mathbf{X} = \operatorname{Spec}(\mathcal{C}(\mathbf{Y};\mathbf{R}))\) 典范地等同于 \(\beta \mathbf{Y}\) ，即 \(\mathbf{Y}\) 的 Stone-Čech 紧化。这样的空间 \(\mathbf{Y}\) 称为平坦的。

\begin{enumerate}
\def\labelenumi{(\alph{enumi})}
\setcounter{enumi}{1}
\item
  平坦空间的每个子空间都是平坦的。平坦空间的每个完全正则商空间都是平坦的。平坦空间的每个有限乘积都是平坦的。平坦空间的每个和都是平坦空间。
\item
  在一个平坦空间中，每个可数子空间都是闭的且离散的。特别地，每个可数平坦空间都是离散的。
\item
  每个 Weierstrass 型完全正则空间（《一般拓扑学》第 IX 章 § 1 习题 22）（更不用说每个紧空间）只要是平坦的，就是有限的。
\item
  与非平凡超滤子相伴的空间是极端不连通的（《一般拓扑学》第 I 章 § 11 习题 21），但不是平坦的。
\item
  离散空间 \(\mathbf{N}\) 是平坦的，但它的 Stone-Čech 紧化不是平坦的，因而环 \(\mathcal{C}(\mathbf{N};\mathbf{R})\) 与 \(\mathcal{C}^{\infty}(\mathbf{N};\mathbf{R})\) 的谱不同构。
\end{enumerate}

\begin{enumerate}
\def\labelenumi{\arabic{enumi}.}
\setcounter{enumi}{18}
\tightlist
\item
  \begin{enumerate}
  \def\labelenumii{(\alph{enumii})}
  \tightlist
  \item
    设 \(\phi: A \to B\) 是一个环同态。证明：设 \(\mathfrak{b}\) 是 \(B\) 的任意一个理想，则有 \(\overline{^a\phi(V(b))} = V(\overline{\phi}^{-1}(b))\) 。
  \end{enumerate}
\end{enumerate}

\begin{enumerate}
\def\labelenumi{(\alph{enumi})}
\setcounter{enumi}{1}
\item
  由 (a) 推出：为使 \(^{a}\phi(\text{Spec}(B))\) 在 \(\text{Spec}(A)\) 中稠密，必须且只需 \(\text{Ker}(\phi)\) 是 A 的一个诣零理想。
\item
  证明：若每个 \(b \in \mathbf{B}\) 都可写成 \(b = h\phi(a)\) ，其中 \(h\) 在 \(\mathbf{B}\) 中可逆，而 \(a \in \mathbf{A}\) ，则 \(^a\phi\) 是 \(\operatorname{Spec}(\mathbf{B})\) 到 \(\operatorname{Spec}(\mathbf{A})\) 的一个子空间上的同胚。
\end{enumerate}

\begin{enumerate}
\def\labelenumi{\arabic{enumi}.}
\setcounter{enumi}{19}
\item
  给出一个环同态 \(\phi: A \to B\) 的例子，使得 \(A\) 的每个极大理想都形如 \(\bar{\phi}(m)\) ，其中 \(m\) 是 \(B\) 的一个极大理想，但 \(A\) 的某些素理想不是在 \(\phi\) 下 \(B\) 的理想的原像（取 \(B\) 为一个域）。
\item
  设 \((\mathbf{A}_{\alpha}, \phi_{\beta \alpha})\) 是环的一个正向系，其指标集是有向的，\(\mathbf{A} = \varinjlim \mathbf{A}_{\alpha}\) ，而 \(\phi_{\alpha}: \mathbf{A}_{\alpha} \to \mathbf{A}\) 是典范同态。若记 \(\mathbf{X}_{\alpha} = \operatorname{Spec}(\mathbf{A}_{\alpha})\) ，则 \((\mathbf{X}_{\alpha}, {}^{a}\phi_{\beta \alpha})\) 是拓扑空间的一个逆向系，且若 \(\mathbf{X} = \operatorname{Spec}(\mathbf{A})\) ，则诸 \({}^{a}\phi_{\alpha}: \mathbf{X} \to \mathbf{X}_{\alpha}\) 组成连续映射的一个逆向系。证明 \(u = \varprojlim {}^{a}\phi_{\alpha}\) 是 \(\mathbf{X}\) 到 \(\varprojlim {}^{\infty}\mathbf{X}_{\alpha}\) 上的一个同胚。（首先证明
\end{enumerate}

：若对所有 \(\alpha\) ，\(p_{\alpha}\) 是 \(\mathbf{A}_{\alpha}\) 的一个素理想，使得 \(p_{\alpha} = \phi_{\beta \alpha}^{-1}(p_{\beta})\) 对所有 \(\alpha \leqslant \beta\) 成立，则 \(p\) （诸 \(\phi_{\alpha}(p_{\alpha})\) 的并）是 \(\mathbf{A}\) 的一个素理想；反之，设 \(p\) 是 \(\mathbf{A}\) 的任意素理想，则它是诸 \(\phi_{\alpha}(p_{\alpha})\) 的并，其中 \(p_{\alpha} = \overline{\phi}_{\alpha}(p)\) 。最后，若 \(f_{\alpha} \in \mathbf{A}_{\alpha}\) 且 \(f = \phi_{\alpha}(f_{\alpha})\) ，则注意：关系 \(p \in \mathbf{X}_f\) 等价于

\[
u (\mathfrak {p}) \in \overline {{\operatorname{pr}}} _ {\alpha} ^ {- 1} ((\mathrm{X} _ {\alpha}) _ {f _ {\alpha}}).
\]

\begin{enumerate}
\def\labelenumi{\arabic{enumi}.}
\setcounter{enumi}{21}
\tightlist
\item
  \begin{enumerate}
  \def\labelenumii{(\alph{enumii})}
  \tightlist
  \item
    设 M 是一个 A-模，\((\mathbf{N}_{i})\) 是 M 的子模的一个有限族。证明 \(\operatorname{Supp}\left(\mathbf{M}/\bigcap_{i}\mathbf{N}_{i}\right)=\bigcup_{i}\operatorname{Supp}(\mathbf{M}/\mathbf{N}_{i})\) （注意 \(M/\bigcap_{i}N_{i}\) 同构于诸 \(M/N_{i}\) 的直和的一个子模）。
  \end{enumerate}
\end{enumerate}

\begin{enumerate}
\def\labelenumi{(\alph{enumi})}
\setcounter{enumi}{1}
\item
  设 \(p\) 是一个素数；\(\mathbf{Z}\) -模 \(\mathbf{Z}\) 的支撑集不同于诸 \(\mathbf{Z}\) -模 \(\mathbf{Z} / p^k\mathbf{Z}\) （ \(k \in \mathbf{N}\) ）的支撑集的并的闭包。
\item
  设 M 是诸 Z-模 \(Z/p^{k}Z\) （ \(k \in N\) ）的直和；证明 Supp(M) 是闭的，但不同于 V(a)，其中 a 是 M 的零化子。
\item
  设 \(\mathbf{N}\) 是诸 \(\mathbf{Z}\) -模 \(\mathbf{Z} / n\mathbf{Z}\) （ \(n \geqslant 1\) ）的直和；证明 \(\operatorname{Supp}(\mathbf{N})\) 在 \(\operatorname{Spec}(\mathbf{Z})\) 中不是闭的。
\item
  由 (c) 与 (d) 推出一个 A-模 M 的例子，使得：若 a 是它的零化子，则 Supp(M) 不是闭的，且其闭包不同于 V(a)。
\item
  证明 \(\mathbf{Z}\) -模 \(\prod_{k=1}^{n} (\mathbf{Z}/p^k\mathbf{Z})\) 的支撑集不同于诸因子模 \(\mathbf{Z}/p^k\mathbf{Z}\) 的支撑集的并的闭包。
\end{enumerate}

\begin{enumerate}
\def\labelenumi{\arabic{enumi}.}
\setcounter{enumi}{22}
\item
  设 \(p\) 是一个素数；给出 \(\mathbf{Z}\) -模 \(M, N\) 的一个例子，其中 \(N\) 是有限生成的，使得 \(\mathbf{M}_{(p)} \neq 0\) 且 \(\mathbf{N}_{(p)} \neq 0\) ，但 \((\mathbf{M} \otimes_{\mathbf{Z}} \mathbf{N})_{(p)} = 0\) （取 \(\mathbf{M} = \mathbf{Q}\) ）。
\item
  \begin{enumerate}
  \def\labelenumii{(\alph{enumii})}
  \tightlist
  \item
    设 A 是一个环，M、N 是两个 A-模，其中 M 是有限生成的。证明 \(\operatorname{Supp}(\operatorname{Hom}_{\mathbf{A}}(\mathbf{M}, \mathbf{N}))\) 包含于交 \(\operatorname{Supp}(\mathbf{M}) \cap \operatorname{Supp}(\mathbf{N})\) 。
  \end{enumerate}
\end{enumerate}

\begin{enumerate}
\def\labelenumi{(\alph{enumi})}
\setcounter{enumi}{1}
\item
  给出一个例子，其中 M 与 N 都是有限表示的，而 \(\operatorname{Supp}(\operatorname{Hom}_{\mathbf{A}}(\mathbf{M},\mathbf{N}))\) 真包含于 \(\operatorname{Supp}(\mathbf{M}) \cap \operatorname{Supp}(\mathbf{N})\) （取 \(\mathbf{A} = \mathbf{Z}_{(p)}\) ，\(\mathbf{N} = \mathbf{A}\) ）。
\item
  设 M 是诸 \(Z/p^{k}Z\) （p 素，\(k \in N\) ）的直和组成的 Z-模。证明 \(\text{Supp}(\text{Hom}_{\mathbf{Z}}(\text{M}, \text{M}))\) 不包含于 \(\text{Supp}(\text{M})\) 。
\end{enumerate}

\begin{enumerate}
\def\labelenumi{\arabic{enumi}.}
\setcounter{enumi}{24}
\tightlist
\item
  设 A 是一个 Noether 环，\(X = \text{Spec}(A)\) 是它的谱，M 是一个有限生成 A-模，\(Y = \text{Supp}(M)\) 。设 \(Y_k (1 \leqslant k \leqslant n)\) 是 Y 的互不相同的连通分支。
\end{enumerate}

\begin{enumerate}
\def\labelenumi{(\alph{enumi})}
\item
  证明存在 M 作为直和的唯一分解 \(M = M_{1} \oplus M_{2} \oplus \cdots \oplus M_{n}\) ，使得对所有 k 有 \(\text{Supp}(M_{k}) = Y_{k}\) 。
\item
  若 \(\mathfrak{a} = \operatorname{Ann}(\mathbf{M})\) ，\(\mathfrak{a}_k = \operatorname{Ann}(\mathbf{M}_k)\) ，证明诸 \(\mathfrak{a}_k\) 两两互素，且 \(\prod_{k=1}^{n} \mathfrak{a}_k = \mathfrak{a}\) 。
\end{enumerate}

（利用第 4 节命题 17，归结到 \(\mathfrak{a} = 0\) 且 \(\mathbf{Y} = \mathbf{X}\) 的情形，并应用第 3 节命题 15。）

\begin{enumerate}
\def\labelenumi{\arabic{enumi}.}
\setcounter{enumi}{25}
\tightlist
\item
  设 \(\phi: \mathbf{A} \to \mathbf{B}\) 是一个环同态，使得映射
\end{enumerate}

\[
^ {a} \phi : \operatorname{Spec} (\mathbf {B}) \rightarrow \operatorname{Spec} (\mathbf {A})
\]

是满射。设 N 是一个有限生成 A-模；证明：若 \(u: M \to N\) 是一个 A-模同态，使得 \(u \otimes 1: M_{(B)} \to N_{(B)}\) 是满射，则 u 是满射（对 Coker(u) 应用第 4 节命题 19）。

\begin{enumerate}
\def\labelenumi{\arabic{enumi}.}
\setcounter{enumi}{26}
\item
  给出一个局部同态 \(\rho: \mathbf{A} \to \mathbf{B}\) 和一个 A-模 \(\mathbf{M}\) （不是有限生成的）的例子，使得 \(\operatorname{Supp}(\mathbf{M}_{(\mathbf{B})})\) 不等于 \(a\rho^{-1}(\operatorname{Supp}(\mathbf{M}))\) （取 \(\mathbf{A} = \mathbf{Z}_{(p)}\) 和 \(\mathbf{B} = \mathbf{Z}/p\mathbf{Z}\) ，这里 \(p\) 是某个素数）。
\item
  给出一个 Z-模 M 的例子，使得：对某个满足 \((p) \in \text{Supp}(M)\) 的素数 p，不存在从 M 到 Z/pZ 的 \(\neq 0\) 的 Z-同态。
\end{enumerate}

